% \newcommand*{\orcid}{}
%

\makeatletter
\let\thetitle\@title
\let\theauthor\@author
\makeatother

\newcommand{\togglepaper}[1][0]{
   \bibliography{../localbibliography}
   \papernote{\scriptsize\normalfont
     \theauthor.
     \titleTemp.
     To appear in:
     E. Di Tor \& Herr Rausgeberin (ed.).
     Booktitle in localcommands.tex.
     Berlin: Language Science Press. [preliminary page numbering]
   }
   \pagenumbering{roman}
   \setcounter{chapter}{#1}
   \addtocounter{chapter}{-1}
}
%
\newbool{bookcompile}
\booltrue{bookcompile}
\newcommand{\bookorchapter}[2]{\ifbool{bookcompile}{#1}{#2}}
%
% %\newcommand{\orcid}[1]{}
\graphicspath{ {./figures/} }
%
\let\eachwordone=\itshape


% Seongha
% \newcommand{\koreex}[1]{{\fontspec{Baekmuk Batang}{#1}}}

% Kutida
\newfontfamily\thaifont{Noto Sans Thai} 
\newcommand{\textthai}[1]{{\thaifont #1}} % für Thai


\newfontfamily\chinesefont{Noto Sans CJK SC}
\newcommand{\textchinese}[1]{{\chinesefont #1}} % Für Chinesisch

% wiemer
\def\hyph{-\penalty0\hskip0pt\relax}
% \definechangesauthor[color=NavyBlue]{SH}



\newcolumntype{K}{>{\raggedright\arraybackslash}p{0.6\textwidth}}  % Right column with 0.7 width


% \newfontfamily\krn[Scale=MatchLowercase,BoldFont=SourceHanSerif-Bold.otf]{SourceHanSerif-Regular.otf}
% \AdditionalFontImprint{Source Han Serif KO}
% \XeTeXlinebreaklocale 'ko'
% \renewcommand{\krn}{}



% Rhee
\newfontfamily\koreanfont{Noto Sans CJK KR} 
\newcommand{\textkorean}[1]{{\koreanfont #1}} % für Koreanisch

\renewcommand{\REF}[2][]{(\ref{#2}#1)}